\section{Proof of main results}\label{Section2}

\begin{proof}
[Proof of Proposition~\ref{T1}]
For $i=1,\ldots, n-1$, we def\/ine $\Theta(x_i)$, $\Theta(x_n)$, and $\Theta(g_i)$
by~\eqref{theta1} and~\eqref{theta3}.
It follows from a~straightforward verif\/ication that~$\Theta$ is a~well-def\/ined homomorphism.



It remains to see that~$\Theta$ is injective.
Observe that $\mathsf{H}$ is spanned by the monomials $x^p g$ for $p=(p_1,\dots,p_n)\in {\mathbb{Z}}^n$ and $g\in G$,
where $p_1,\ldots, p_n\ge 0$.
We call $p_1+\dots+p_n$ the total degree of the monomial $x^p g$.
The image of $x^p g$ under~$\Theta$ is the sum of $y^p g$ with terms of strictly smaller total degrees.
Therefore, if $\alpha\in \mathsf{H}$ is nonzero, we can write it as a~sum $\alpha_0 + \alpha_1 + \cdots$, where
$\alpha_k$ is a~linear combination of monomials~$x^p g$ with total degree~$k$.
If~$k$ is the maximal integer with~$\alpha_k$ nonzero, then $\Theta(\alpha_k)$ is nonzero, and hence $\Theta(\alpha)$ is
also nonzero.
\end{proof}

\begin{Remark}
\label{pbwbasis}
It follows from Proposition~\ref{T1} that the monomials $x_1^{p_1}\cdots x_n^{p_n} g$ for non-negative integers $p_1,
\ldots, p_n$ and $g\in G$ form a~basis for $\mathsf{H}$ (called the PBW basis of $\mathsf{H}$).
This was f\/irst proved in~\cite[Example 2.16]{W} using~\cite[Theorem 2.10]{W}.
\end{Remark}

We have an increasing f\/iltration on $\mathsf{H}$ def\/ined by setting $\deg(x_i)=1$ and $\deg(g)=0$ for all
$i\in\{1,\ldots,n\}$, $g\in G$.
It is immediate from Remark~\ref{pbwbasis} that the natural homomorphism $SV\#_\alpha G \to {\mathrm{gr}} \mathsf{H}$ is
an isomorphism, where ${\mathrm{gr}} \mathsf{H}$ denotes the associated graded algebra of $\mathsf{H}$.

The proof of~\eqref{thetaw} in the following lemma is the key calculation in this paper.

\begin{Lemma}\label{keylemma}\quad
\begin{enumerate}\itemsep=0pt
\item[$(i)$] One has:
\begin{gather}
\Theta\big(x_i^\ell\big) = y_i^\ell - \tau_i^\ell y_{i+1}^{-\ell},
\label{thetaxil}
\\
\Theta\big(x_n^\ell\big) = y_n^\ell - (-1)^{n(\ell-1)} \tau_n^\ell y_1^{-\ell},
\label{thetaxnl}
\end{gather}
for all $i\in \{1,\ldots, n-1\}$.

\item[$(ii)$] One has:
\begin{gather}
\label{thetaw}
\Theta(w) = y_1\cdots y_n + (-1)^n \zeta^{n-2} \tau_1\cdots \tau_n y_1^{-1} \cdots y_n^{-1}.
\end{gather}
\end{enumerate}
\end{Lemma}

\begin{proof}
(i) To prove~\eqref{thetaxil}, we need to show that
\begin{gather}
\label{eqexpand1}
\underbrace{\big(y_i - \zeta \tau_i y_{i+1}^{-1}g_i\big)
\cdots
\big(y_i - \zeta \tau_i y_{i+1}^{-1}g_i\big)}_\ell = y_i^\ell - \tau_i^\ell y_{i+1}^{-\ell}.
\end{gather}
Since $g_iy_i=\zeta y_i g_i$ and $g_i y_{i+1}^{-1} = \zeta y_{i+1}^{-1} g_i$, the product on the left hand side
of~\eqref{eqexpand1} is a~linear combination of $y_i^k y_{i+1}^{k-\ell} g_i^{\ell-k}$ for $k=0, 1,\ldots, \ell$.
Moreover, the coef\/f\/icient of $y_i^k y_{i+1}^{k-\ell} g_i^{\ell-k}$ in this linear combination is the same as the
coef\/f\/icient of $u^k$ when we expand the product
\begin{gather}
\label{poly1}
\big(u- \zeta^\ell \tau_i\big) \big(u- \zeta^{\ell-1} \tau_i\big) \cdots (u- \zeta \tau_i)
\end{gather}
in the polynomial ring ${\mathbb{C}}[u]$.
Since the polynomial in~\eqref{poly1} is equal to $u^\ell-\tau_i^\ell$, the identity~\eqref{thetaxil} follows.
The proof of~\eqref{thetaxnl} is similar except that
\begin{gather*}
\underbrace{g_n * \dots * g_n}_\ell = (-1)^{n(\ell-1)}.
\end{gather*}
(ii) For any $h_*=\{h_1<\dots<h_j\} \in I$, we let
\begin{gather*}
h_*' = \{h_r\in h_* \mid h_s-h_r \in \{1, 1-n\}\ \text{for some}\ s \},
\\
\chi(h_*) = |\{h_r\in h'_* \mid h_r\neq n \}| - |\{h_r\in h'_* \mid h_r=n \}|,
\\
E(h_*) = \zeta^{\chi (h_*)} \tau_{h_1} \cdots \tau_{h_j} y^{\delta - \varepsilon_{h_1} - \dots - \varepsilon_{h_j}}
g_{h_1} * \cdots * g_{h_j}.
\end{gather*}
Now suppose $i_*=\{i_1<\dots<i_k\}\in J$.
Let~$D$ be the subset of $\{1,\dots,n\}$ consisting of all~$d$ such that $d \not\equiv i_r,\ i_r+1$ (mod~$n$) for all~$r$.
We denote by $d_1 < \dots < d_p$ the elements of~$D$.
Then
\begin{gather*}
\Theta\big(\tau_{i_1}\cdots \tau_{i_k} x^{\delta-\varepsilon_{i_1} - \dots - \varepsilon_{i_k}} g_{i_1}*\cdots * g_{i_k}\big)
\\
\qquad
= \tau_{i_1}\cdots \tau_{i_k} \left(y_{d_1} - \frac{\zeta t_{d_1}}{\zeta-1}y_{d_1+1}^{-1}g_{d_1} \right) \cdots
\left(y_{d_p} - \frac{\zeta t_{d_p}}{\zeta-1}y_{d_p+1}^{-1}g_{d_p} \right) g_{i_1}*\cdots * g_{i_k}
\\
\qquad
= \tau_{i_1}\cdots \tau_{i_k} \sum\limits_{S\subset D} Y_{d_1}(S) \cdots Y_{d_p}(S) g_{i_1}*\cdots * g_{i_k},
\end{gather*}
where, for $r=1,\ldots, p$,
\begin{gather*}
Y_{d_r}(S) =
\begin{cases}
y_{d_r}, & \text{if} \quad d_r\notin S,
\\
- \zeta(\zeta-1)^{-1} t_{d_r}y_{d_r+1}^{-1}g_{d_r}, & \text{if} \quad d_r\in S.
\end{cases}
\end{gather*}
Setting $h_* = i_* \cup S$, we obtain\footnote{Note that if $d_r\in S$ but $d_r+1\in D-S$, then the term $g_{d_r}$ in
$Y_{d_r}(S)$ appears on the left of the term $y_{d_r+1}$ of $Y_{d_r+1}(S)$ and one has $g_{d_r} y_{d_r+1} = \zeta^{-1}
y_{d_r+1} g_{d_r}$.
However, if $n\in S$ but $1\in D-S$, then the term $g_n$ in $Y_n(S)$ already appears to the right of the term $y_1$ of
$Y_1(S)$.
This is the reason why the def\/inition of $\tau_n$ dif\/fers from the corresponding def\/initions of $\tau_1,\ldots,
\tau_{n-1}$ by a~factor of~$\zeta$.}
\begin{gather*}
\Theta\big(\tau_{i_1}\cdots \tau_{i_k} x^{\delta-\varepsilon_{i_1} - \dots - \varepsilon_{i_k}} g_{i_1}*\cdots * g_{i_k}\big) =
\sum\limits_{\{h_*\in I \mid i_*\subset h_*- h'_* \}} (-1)^{|h_*|-|i_*|} E(h_*).
\end{gather*}
Hence,
\begin{gather*}
\Theta(w) = \sum\limits_{\{i_1<\dots<i_k\}\in J} \Theta\big(\tau_{i_1}\cdots \tau_{i_k} x^{\delta-\varepsilon_{i_1} - \dots -
\varepsilon_{i_k}} g_{i_1}*\dots * g_{i_k}\big)
\\
\hphantom{\Theta(w)}{}
{} = \sum\limits_{i_*\in J} \left(\sum\limits_{\{h_*\in I \mid i_*\subset h_*- h'_*\}} (-1)^{|h_*|-|i_*|} E(h_*) \right)
 = \sum\limits_{h_*\in I} \left (E(h_*) \sum\limits_{i_*\subset h_*-h'_*} (-1)^{|h_*|-|i_*|} \right).
\end{gather*}
If $|h_*|=n$, then $h'_*= h_*$.
If $|h_*|\notin\{0,n\}$, then $h'_* \neq h_*$.
Therefore,
\begin{gather*}
E(h_*) \sum\limits_{i_*\subset h_*-h'_*} (-1)^{|h_*|-|i_*|} =
\begin{cases}
y_1\cdots y_n & \text{if} \quad |h_*|=0,
\\
(-1)^n \zeta^{n-2} \tau_1\cdots \tau_n y_1^{-1} \cdots y_n^{-1} & \text{if} \quad |h_*|=n,
\\
0 & \text{else.}
\end{cases}
\tag*{\qed}
\end{gather*}
\renewcommand{\qed}{}
\end{proof}

\begin{proof}[Proof of Theorem~\ref{T2}] It is easy to see that the center of $SV\#_\alpha G$ is the algebra of~$G$-invariant
elements $(SV)^G$ of $SV$, and moreover, the algebra $(SV)^G$ is generated by $x_i^\ell$ ($i=1,\ldots,n$) and $x_1\cdots x_n$.

Using Lemma~\ref{keylemma}, we see that
\begin{gather*}
\Theta\big(x_i^\ell\big)
\quad
\text{for}\quad i=1,\ldots, n,
\qquad
\text{and}
\qquad
\Theta(w)
\end{gather*}
are in the center of ${\mathbb{C}}[y_1^\pm, \ldots, y_n^\pm]\#_\alpha G$.
Since the homomorphism~$\Theta$ is injective, the ele\-ments~$x_i^\ell$ ($i=1,\ldots, n$) and~$w$ are in the center of
$\mathsf{H}$.
Since the principal symbols of $x_1^\ell, \ldots, x_n^\ell$ and~$w$ in $SV\#_\alpha G$ are, respectively, $x_1^\ell,
\ldots, x_n^\ell$ and $x_1\cdots x_n$, the theorem follows from a~standard argument.
\end{proof}
